% Generated draft: requires mathematical and source review.
\documentclass[11pt]{article}
\usepackage[margin=25mm]{geometry}
\usepackage{fontspec}
\setmainfont{Latin Modern Roman}
\usepackage{amsmath,amssymb,mathtools,unicode-math}
\setmathfont{Latin Modern Math}
\usepackage{enumitem,needspace,xurl,hyperref}
\setlength{\parindent}{0pt}
\setlength{\parskip}{5pt}
\setlength{\emergencystretch}{3em}
\newcommand{\N}{\mathbb{N}}
\newcommand{\Z}{\mathbb{Z}}
\newcommand{\Q}{\mathbb{Q}}
\newcommand{\R}{\mathbb{R}}
\newcommand{\C}{\mathbb{C}}
\newcommand{\F}{\mathbb{F}}
\newcommand{\A}{\mathbb{A}}
\newcommand{\PP}{\mathbb{P}}
\newcommand{\E}{\mathbb{E}}
\newcommand{\eps}{\varepsilon}
\newcommand{\dd}{\,\mathrm{d}}
\begin{document}
\section*{UnsolvedMath ID: 30003322}
\textbf{Source problem code:} catalogue entry
\subsection{Definitions and mathematical statement}
Let \(\mathbf{No}\) denote the proper class of surreal numbers with its usual ordered-field operations and simplicity relation \(\prec\). In the sign-sequence representation, \(y\prec x\) means that the sign sequence of \(y\) is a proper initial segment of that of \(x\). A subclass \(A\subseteq\mathbf{No}\) is \emph{initial} if
\[
 x\in A\ \text{and}\ y\prec x\quad\Longrightarrow\quad y\in A.
\]
An initial subgroup is an initial subclass closed under the additive-group operations, with the inherited order. An initial subfield is an initial subclass that is a subfield with the inherited order. Isomorphisms are understood to preserve the relevant group or field operations and order.

Write \(\mathrm{DensOAG}\) for the theory of nontrivial densely ordered abelian groups: the order is total and translation invariant, the group has a nonzero element, and between any two distinct ordered elements there is another element. Write \(\mathrm{DOAG}\) for the theory of nontrivial divisible ordered abelian groups, where divisibility means
\[
 \forall n\in\mathbb{Z}_{\geq1}\ \forall a\ \exists b\quad nb=a.
\]
Write \(\mathrm{OF}\) for the theory of ordered fields and \(\mathrm{RCF}\) for the theory of real-closed ordered fields. An ordered field is real closed if every positive element has a square root and every polynomial of odd degree has a root.

For a theory \(T\), consider the following properties, quantifying only over models whose universes are sets:
\[
\begin{aligned}
 P_{\mathrm{grp}}(T)\quad &\Longleftrightarrow\quad
 \text{every set-sized }M\models T\text{ is isomorphic to an initial subgroup of }\mathbf{No},\\
 P_{\mathrm{fld}}(T)\quad &\Longleftrightarrow\quad
 \text{every set-sized }M\models T\text{ is isomorphic to an initial subfield of }\mathbf{No}.
\end{aligned}
\]
Here the isomorphic image of each set-sized model is itself a set, although \(\mathbf{No}\) is a proper class.

\textbf{Requested foundational statement.} Can the only-if directions of the source report's Theorem~6 be proved in von Neumann--Bernays--G\"odel set theory (NBG), using only models with set-sized universes? The supplied reformulation describes these directions schematically as
\[
\begin{aligned}
 T\supseteq\mathrm{DensOAG},\quad P_{\mathrm{grp}}(T)
 &\quad\Longrightarrow\quad T\equiv\mathrm{DOAG},\\
 T\supseteq\mathrm{OF},\quad P_{\mathrm{fld}}(T)
 &\quad\Longrightarrow\quad T\equiv\mathrm{RCF}.
\end{aligned}
\]
In these displays, \(T\equiv S\) provisionally means that \(T\) and \(S\) have exactly the same set-sized models in the relevant common language; it does not mean literal equality of axiom lists.

\textbf{Unresolved qualifications.} The supplied record does not reproduce Theorem~6 or specify its full hypotheses. In particular, it does not state whether theories must be consistent or have a set-sized model, whether they are first-order theories in exactly the usual ordered-group or ordered-field language, or whether additional language symbols are permitted. Without a nonvacuity hypothesis, the condition on every model also applies vacuously to a theory having no set-sized models. The intended meaning of ``the theory of'' must also be checked against the theorem: equality of model classes, equality of deductive closures, and merely requiring every model to be divisible or real closed must not be silently identified without the relevant hypotheses. Thus the displayed implications record the supplied intended target, subject to these unresolved qualifications; they are not asserted as established theorems.
\subsection{Short English statement}
The question asks whether NBG can prove the stated characterizations of theories whose models occur as initial surreal subgroups or subfields, while considering only models with set-sized universes. The exact hypotheses and meaning of equality of theories require checking against the report's Theorem~6. This is a statement-only draft requiring human mathematical review.
\subsection{Sources}
\begin{itemize}
\item Local registry: \texttt{data/unsolvedmath/problems.json}, supplied source record for permanent UnsolvedMath ID \(30003322\), \emph{Set-Sized Models of Initial Surreal Substructures}; exact source problem code \texttt{catalogue entry}. The record belongs to the Oberwolfach Reports Open Problems collection and supplies no individual proposer.
\item Supplied primary-source citation: \emph{Mini-Workshop: Surreal Numbers, Surreal Analysis, Hahn Fields and Derivations} (2017), \emph{Oberwolfach Reports}, pp.~3313--3372, DOI \texttt{10.4171/owr/2016/60}. Original OWR report: \url{https://doi.org/10.4171/owr/2016/60}. The original statement specifically refers to the only-if directions of Theorem~6. The supplied proposed year and year tag are 2016, whereas the supplied citation labels the report 2017; both are retained without resolving their relationship.
\item Supplied literature source: Rangel--Mariano, \url{https://arxiv.org/abs/1911.12726}. No title or specific supporting passage is supplied.
\item Literature assessment checked at \texttt{2026-08-21}: the supplied curation labels the problem open and reports that no later primary result located in that review establishes the requested NBG reduction. The listed checked sources are the original OWR report and Rangel--Mariano, linked above. This is supplied dated curation, not an independently verified literature assessment.
\item The registry labels the statement \texttt{corrected\_verified} and reports that the official passage was checked before its self-contained reformulation. This verification history is supplied metadata; no independent source inspection or mathematical verification is claimed here.
\end{itemize}
\end{document}
